\section*{Capítulo \ref{ch:b2:fourier} --- Séries de Fourier}

\begin{solution}{exo:b2:fourier:1}
A estranheza mata o $a_n $. Para $ n \geq 1$:
\[
b_n = \frac{2}{\pi}\int_0^{\pi} \sin nt\,\dd t
= \frac{2}{\pi}\cdot\frac{1 - (-1)^n}{n}
= \begin{cases} \frac{4}{\pi n} & n \text{ odd},\\ 0 & n
\text{ even}. \end{cases}
\]
Então $ S(f)(t) = \frac{4}{\pi}\sum_{k\geq0} \frac{\sin\bigl((2k+1)t
\bigr)}{2k+1}$. Dirichlet em $ t = \frac\pi2$ (um continuidade ponto,
valor $1$): $\sin\bigl((2k+1)\frac\pi2\bigr) = (-1)^k $, dando
\[
1 = \frac4\pi \sum_{k\geq0}\frac{(-1)^k}{2k+1}
\quad\Longrightarrow\quad
\sum_{k\geq0}\frac{(-1)^k}{2k+1} = \frac{\pi}{4}
\quad\text{(Leibniz)}.
\]
No salto $ t = 0$: a série soma a $0 = \frac{f(0^+) +
f(0^-)}{2}$, como prescreve Dirichlet (todo termo desaparece:
consistente).
\end{solution}

\begin{solution}{exo:b2:fourier:2}
A uniformidade mata o $ b_n $; $ a_0 = \frac{1}{\pi}\int_{-\pi}^\pi\abs
t\,\dd t = \pi $, e para $ n \geq 1$:
\[
a_n = \frac{2}{\pi}\int_0^\pi t\cos nt\,\dd t
= \frac{2}{\pi}\cdot\frac{(-1)^n - 1}{n^2}
= \begin{cases} -\frac{4}{\pi n^2} & n \text{ odd},\\ 0 & n
\text{ even}, \end{cases}
\]
(uma integração por partes). Por isso
\[
\abs t = \frac{\pi}{2} - \frac{4}{\pi}\sum_{k\geq0}
\frac{\cos\bigl((2k+1)t\bigr)}{(2k+1)^2} ,
\]
com \emph{normal} convergência ($\sum (2k+1)^{-2} < \infty $) --- como
\cref{thm:b2:fourier:parseval}~(1) prevê isso \omterm{def:b2:metric:continuity}{contínuo}
por partes-$ C^1$ função. No $ t = 0$:
\[
0 = \frac\pi2 - \frac4\pi\sum_{k\geq0}\frac{1}{(2k+1)^2}
\quad\Longrightarrow\quad
\sum_{k\geq0}\frac{1}{(2k+1)^2} = \frac{\pi^2}{8} .
\]
Divisão $\sum \frac{1}{n^2}$ em partes pares e ímpares: $ S =
\frac{\pi^2}{8} + \frac S4$, então $ S = \frac{\pi^2}{6}$: Basileia novamente.
\end{solution}

\begin{solution}{exo:b2:fourier:3}
Uniformidade: $ b_n = 0$; $ a_0 = \frac{1}{\pi}\int_{-\pi}^{\pi} t^2 =
\frac{2\pi^2}{3}$; duas integrações por partes dão $ a_n =
\frac{4(-1)^n}{n^2}$ ($ n \geq 1$). Por isso
\[
t^2 = \frac{\pi^2}{3} + 4\sum_{n\geq1}
\frac{(-1)^n}{n^2}\cos nt
\qquad (\abs t \leq \pi),
\]
\omterm{def:b2:funcseq:series}{normalmente} convergente. No $ t = 0$: $0 = \frac{\pi^2}{3} +
4\sum\frac{(-1)^n}{n^2}$, ou seja\ $\sum_{n\geq1}
\frac{(-1)^{n+1}}{n^2} = \frac{\pi^2}{12}$. Parseval:
\[
\frac{1}{2\pi}\int_{-\pi}^{\pi} t^4\,\dd t = \frac{\pi^4}{5}
= \frac{a_0^2}{4} + \frac12\sum_{n\geq1} a_n^2
= \frac{\pi^4}{9} + 8\sum_{n\geq1}\frac{1}{n^4} ,
\]
então $\sum \frac{1}{n^4} = \frac18\bigl(\frac{\pi^4}{5} -
\frac{\pi^4}{9}\bigr) = \frac{\pi^4}{90}$.
\end{solution}

\begin{solution}{exo:b2:fourier:4}
Se além disso $ f $ é por partes $ C^1$: Parseval (provado) dá
$\norm f_2^2 = \sum\abs{c_n}^2 = 0$ e estrita positividade do
integral do \omterm{def:b2:metric:continuity}{contínuo} $\abs f^2$ forças $ f = 0$.

Por apenas \omterm{def:b2:metric:continuity}{contínuo} $ f $, admita Parseval geral: mesmo unifilar
prova. (Esboço da rota de densidade: Fejtipo er/Weierstrass
aproximação trigonométrica mostra que polinômios trigonométricos são
$\norm\cdot_2$-denso entre \omterm{def:b2:metric:continuity}{contínuo} funções periódicas; desde $ f
\perp $ todos eles, $\norm f_2^2 = \langle f, f - P\rangle \leq
\norm f_2\norm{f - P}_2$ para aproximantes $ P $, forçando $\norm f_2 =
0$.)
\end{solution}

\begin{solution}{exo:b2:fourier:5}
Uniformidade: $ b_n = 0$;
\[
a_n = \frac{2}{\pi}\int_0^\pi \cos(\alpha t)\cos(nt)\,\dd t
= \frac{2}{\pi}\cdot
\frac{(-1)^n\,\alpha\sin(\pi\alpha)}{\alpha^2 - n^2}
\]
(produto para soma, depois integre; $ a_0 =
\frac{2\sin(\pi\alpha)}{\pi\alpha}$). Dirichlet em $ t = \pi $ (um
continuidade ponto da extensão periódica, cujos valores unilaterais
concordar por igualdade):
\[
\cos(\pi\alpha)
= \frac{\sin(\pi\alpha)}{\pi\alpha}
+ \sum_{n\geq1} \frac{2\alpha\sin(\pi\alpha)}{\pi(\alpha^2 -
n^2)}\,(-1)^n\cos(n\pi)
= \frac{\sin(\pi\alpha)}{\pi}\Bigl(\frac{1}{\alpha} +
\sum_{n\geq1}\frac{2\alpha}{\alpha^2 - n^2}\Bigr),
\]
usando $(-1)^n\cos n\pi = 1$. Dividindo por $\sin(\pi\alpha)/\pi $
dá a expansão cotangente.
\end{solution}

\begin{solution}{exo:b2:fourier:6}
Iterando $ c_n(f') = \iu n\,c_n(f)$ ($ k $ vezes, integração por
peças em todo o $ C^{k}$ peças com valores limite correspondentes):
$ c_n(f^{(k)}) = (\iu n)^k c_n(f)$. Os coeficientes do
por partes \omterm{def:b2:metric:continuity}{contínuo} $ f^{(k)}$ são limitados (na verdade $\to 0$, Bessel):
\[
\abs{c_n(f)} = \frac{\abs{c_n(f^{(k)})}}{\abs n^k}
= O\bigl(\abs n^{-k}\bigr) .
\]
\end{solution}

\begin{solution}{exo:b2:fourier:7}
Parseval para $ f $ e para $ f'$ (ambos legítimos: $ f $ é $ C^1$, $ f'$
por partes \omterm{def:b2:metric:continuity}{contínuo} --- de fato \omterm{def:b2:metric:continuity}{contínuo}):
\[
\frac{1}{2\pi}\int \abs f^2 = \sum_{n\neq0} \abs{c_n}^2
\quad (c_0 = 0 \text{ by the mean-zero hypothesis}),
\qquad
\frac{1}{2\pi}\int \abs{f'}^2 = \sum_{n\neq0} n^2\abs{c_n}^2 .
\]
Termo, $ n^2\abs{c_n}^2 \geq \abs{c_n}^2$ para $\abs n \geq 1$:
a desigualdade segue. Forças de igualdade $(n^2 - 1)\abs{c_n}^2 = 0$
para todos $ n $, ou seja\ $ c_n = 0$ para $\abs n \geq 2$: $ f(t) =
c_1\eu^{\iu t} + c_{-1}\eu^{-\iu t} = a\cos t + b\sin t $ (real
formulário); inversamente tal $ f $ dê igualdade.
\end{solution}

\begin{solution}{exo:b2:fourier:8}
De \cref{exo:b2:fourier:1}, $ S_{2N-1}(x) =
\frac4\pi\sum_{k=0}^{N-1}\frac{\sin((2k+1)x)}{2k+1}$. No $ x_N =
\frac{\pi}{2N}$, com $ u_k = (2k+1)x_N = \frac{(2k+1)\pi}{2N}$:
\[
S_{2N-1}(x_N) = \frac{4}{\pi}\sum_{k=0}^{N-1}\frac{\sin u_k}{2k+1}
= \frac{4}{\pi}\sum_{k=0}^{N-1}\frac{\sin u_k}{u_k}\cdot
\frac{u_k}{2k+1}
= \frac{2}{\pi}\sum_{k=0}^{N-1}\frac{\sin u_k}{u_k}\cdot
\frac{\pi}{N},
\]
desde $\frac{u_k}{2k+1} = \frac{\pi}{2N}$. Os pontos $ u_k $ são os
pontos médios do $ N $ subintervalos de $\intcc{0}{\pi}$ de comprimento
$\frac{\pi}{N}$: a soma é uma soma de Riemann do ponto médio do
\omterm{def:b2:metric:continuity}{contínuo} $ u \mapsto \frac{\sin u}{u}$, portanto converge para
\[
\frac{2}{\pi}\int_0^\pi \frac{\sin u}{u}\,\dd u
\approx \frac{2}{\pi}\times 1.8519 \approx 1.179 .
\]
As somas parciais próximas ao salto ultrapassam o valor $1$ por $\approx
18\%$ do meio salto para sempre: o fenômeno de Gibbs, quantificado.
\end{solution}

\begin{solution}{exo:b2:fourier:9}
De $\cos 3t = 4\cos^3t - 3\cos t $:
\[
\cos^3 t = \frac{3\cos t + \cos 3t}{4},
\qquad
\sin^2t\,\cos t = \cos t - \cos^3 t
= \frac{\cos t - \cos 3t}{4} .
\]
Cada um é um polinômio trigonométrico, portanto igual ao seu próprio
Série de Fourier (unicidade dos coeficientes: duas expansões
diferiria por um polinômio trigonométrico com todos
coeficientes zero). Para $\cos^3t $: $ a_1 = \frac34$, $ a_3 =
\frac14$, todos os outros $ a_n $ e tudo $ b_n $ zero; $ c_{\pm1} =
\frac38$, $ c_{\pm3} = \frac18$. Para $\sin^2t\cos t $: $ a_1 =
\frac14$, $ a_3 = -\frac14$; $ c_{\pm1} = \frac18$, $ c_{\pm3} =
-\frac18$.
\end{solution}

\begin{solution}{exo:b2:fourier:10}
Cálculo direto:
\[
c_n = \frac{1}{2\pi}\int_{-\pi}^{\pi}\eu^{(a - \iu n)t}\dd t
= \frac{\eu^{(a-\iu n)\pi} - \eu^{-(a - \iu n)\pi}}
{2\pi(a - \iu n)}
= \frac{(-1)^n\sinh(a\pi)}{\pi(a - \iu n)} ,
\]
usando $\eu^{\pm\iu n\pi} = (-1)^n $. No $ t = \pi $ o periódico
extensão salta de $\eu^{a\pi}$ a $\eu^{-a\pi}$; Dirichlet
(\omterm{def:b2:quadratic:adjoint}{simétrico} somas parciais) dá
\[
\cosh(a\pi) = \sum_{n\in\Z}(-1)^n c_n\,
= \frac{\sinh(a\pi)}{\pi}\Bigl(\frac1a +
\sum_{n\geq1}\Bigl(\frac{1}{a - \iu n} +
\frac{1}{a + \iu n}\Bigr)\Bigr)
= \frac{\sinh(a\pi)}{\pi}\Bigl(\frac1a +
\sum_{n\geq1}\frac{2a}{a^2 + n^2}\Bigr) ,
\]
as partes imaginárias dos termos emparelhados cancelando. Dividir por
$\sinh(a\pi)$:
\[
\coth(\pi a) = \frac{1}{\pi a} +
\sum_{n\geq1}\frac{2a}{\pi(a^2 + n^2)} ,
\]
o gêmeo hiperbólico de \cref{exo:b2:fourier:5}.
\end{solution}

\begin{solution}{exo:b2:fourier:11}
Comutatividade: substituto $ s = x - t $ e usar periodicidade de
o integrando. Para $ c_n(f * g)$, o integrando $(x, t) \mapsto
f(x-t)g(t)\eu^{-\iu nx}$ é \omterm{def:b2:metric:continuity}{contínuo}; os dois iteraram
integrais concordam (ambos, como funções do limite superior do
variável externa, desaparece no extremo esquerdo e tem o mesmo
derivada --- continuidade mais
\cref{thm:b2:integration:continuity} justificar a diferenciação
a integral iterada). Por isso
\[
c_n(f*g) = \frac{1}{2\pi}\int_{-\pi}^{\pi} g(t)\,\eu^{-\iu nt}
\Bigl(\frac{1}{2\pi}\int_{-\pi}^{\pi}
f(x-t)\,\eu^{-\iu n(x-t)}\dd x\Bigr)\dd t
= c_n(f)\,c_n(g),
\]
o ser integral interno $ c_n(f)$ para cada $ t $ (substituição
e periodicidade). Finalmente \cref{lem:b2:fourier:kernel} diz
$ S_N(f)(x) = \frac{1}{2\pi}\int f(x+u)D_N(u)\dd u $; o
substituição $ u \mapsto -t $ e a uniformidade de $ D_N $ vire isso
em $(f * D_N)(x)$.
\end{solution}

\begin{solution}{exo:b2:fourier:12}
\begin{enumerate}
  \item Com $ w = r\eu^{\iu t}$:
        \[
        P_r(t) = 1 + 2\,\Re\frac{w}{1 - w}
        = \Re\frac{1 + w}{1 - w}
        = \frac{1 - \abs w^2}{\abs{1 - w}^2}
        = \frac{1 - r^2}{1 - 2r\cos t + r^2} > 0 .
        \]
        Significar $1$: integração termo a termo do \omterm{def:b2:funcseq:series}{normalmente}
        série convergente mantém apenas $ n = 0$.
  \item Para $\delta \leq \abs t \leq \pi $: $\cos t \leq
        \cos\delta $, então $ P_r(t) \leq \frac{1 - r^2}{1 -
        2r\cos\delta + r^2}$, cujo denominador tende a $2 -
        2\cos\delta > 0$ enquanto o numerador tende a $0$:
        convergência uniforme para $0$ fora de qualquer bairro de
        $0$.
  \item Integração termo a termo (convergência normal em $ t $)
        dá $(f * P_r)(x) = \sum_n r^{\abs
        n}c_n(f)\eu^{\iu nx}$. A identidade aproximada
        argumento: com média $1$ e positividade,
        \[
        \abs{(f*P_r)(x) - f(x)}
        \leq \frac{1}{2\pi}\int_{-\pi}^{\pi}
        \abs{f(x-t) - f(x)}\,P_r(t)\,\dd t ,
        \]
        dividir em $\abs t = \delta $: no máximo $\varepsilon $
        (Heine) mais $2\norm f_\infty\sup_{\delta\leq\abs
        t\leq\pi}P_r \to \varepsilon $: convergência uniforme como
        $ r \to 1^-$. Isso é \omterm{thm:b2:series:abel}{Resumo de Abel} do Fourier
        série --- o gêmeo Fourier da série de potência
        teoria dos limites do capítulo.
\end{enumerate}
\end{solution}

\begin{solution}{pb:b2:fourier:1}
\textbf{1.} Média \cref{lem:b2:fourier:kernel} sobre $ n = 0,
\dots, N-1$ (linearidade da integral) dá $\sigma_N(f)(x) =
\frac{1}{2\pi}\int f(x+u)F_N(u)\dd u $; cada $ D_n $ tem significado
$1$, então $ F_N $ tem significado $1$.

\textbf{2.} Com $\eu^{\iu u} - 1 = 2\iu\,\eu^{\iu
u/2}\sin\frac u2$:
\[
\sum_{n=0}^{N-1}\sin\Bigl(\Bigl(n + \frac12\Bigr)u\Bigr)
= \Im\Bigl[\eu^{\iu u/2}\,\frac{\eu^{\iu Nu} - 1}{\eu^{\iu u}
- 1}\Bigr]
= \Re\,\frac{1 - \eu^{\iu Nu}}{2\sin\frac u2}
= \frac{1 - \cos Nu}{2\sin\frac u2}
= \frac{\sin^2\frac{Nu}2}{\sin\frac u2} .
\]
Dividindo por $ N\sin\frac u2$:
\[
F_N(u) = \frac1N\sum_{n=0}^{N-1}
\frac{\sin\bigl((n+\frac12)u\bigr)}{\sin\frac u2}
= \frac{1}{N}\,
\frac{\sin^2\frac{Nu}{2}}{\sin^2\frac u2} \geq 0 .
\]

\textbf{3.} Sobre $\delta \leq \abs u \leq \pi $: $\sin^2\frac u2
\geq \sin^2\frac\delta2$ e $\sin^2\frac{Nu}2 \leq 1$: $ F_N
\leq \frac{1}{N\sin^2(\delta/2)} \to 0$, \omterm{def:b2:funcseq:def}{uniformemente} lá.

\textbf{4.} Pela média da unidade, $\sigma_Nf(x) - f(x) =
\frac{1}{2\pi}\int\bigl(f(x+u) - f(x)\bigr)F_N(u)\dd u $. Dado
$\varepsilon $, Heine cede $\delta $ com $\abs{f(x+u) - f(x)}
\leq \varepsilon $ para $\abs u \leq \delta $, \omterm{def:b2:funcseq:def}{uniformemente} em $ x $.
Então, usando $ F_N \geq 0$ e sua média unitária,
\[
\abs{\sigma_Nf(x) - f(x)}
\leq \varepsilon +
2\norm f_\infty\cdot\frac{1}{2\pi}
\int_{\delta\leq\abs u\leq\pi}F_N
\leq \varepsilon + \frac{2\norm
f_\infty}{N\sin^2\frac\delta2}
\leq 2\varepsilon
\]
para grande $ N $, \omterm{def:b2:funcseq:def}{uniformemente} em $ x $: Fejteorema de er.

\textbf{5.} $ F_N $ é par com média $1$: cada metade
$\intcc{0}{\pi}$, $\intcc{-\pi}{0}$ carrega maldade $\frac12$.
Então
\[
\sigma_Nf(x) - \frac{f(x^+)+f(x^-)}{2}
= \frac{1}{2\pi}\int_0^\pi\bigl(f(x+u) -
f(x^+)\bigr)F_N\,\dd u
+ \frac{1}{2\pi}\int_{-\pi}^0\bigl(f(x+u) -
f(x^-)\bigr)F_N\,\dd u ;
\]
em cada metade, dividido em $\abs u = \delta $ onde o unilateral
limite é $\varepsilon $-feche e deixe a pergunta 3 matar o distante
parte: ambas as integrais tendem a $0$. O limite: $ F_N \geq 0$
dá $\abs{\sigma_Nf(x)} \leq
\frac{1}{2\pi}\int\abs{f(x+u)}F_N \leq \norm f_\infty $.

\textbf{6.} Cada $\sigma_N(f)$ é um polinômio trigonométrico
(uma média do $ S_n(f)$, $ n < N $), e $\sigma_N(f) \to f $
\omterm{def:b2:funcseq:def}{uniformemente}: o teorema trigonométrico de Weierstrass.

\textbf{7.} $ c_n(f) = 0$ para todos $ n $ faz cada $ S_n(f) = 0$,
daí cada $\sigma_N(f) = 0$; por Fejérrrr, $ f = \lim\sigma_Nf =
0$. Aplicando isso a uma diferença: \omterm{def:b2:metric:continuity}{contínuo} periódico
funções são determinadas pela sua \omterm{def:b2:fourier:coefficients}{Coeficientes de Fourier}.

\textbf{8.} $ S_Nf $ é a projeção ortogonal de $ f $ a $\mathcal T_N $ (\cref{prop:b2:fourier:bessel}), então minimiza
$\norm{f - P}_2$ sobre $ P \in \mathcal T_N $; desde $\sigma_Nf
\in \mathcal T_{N-1} \subseteq \mathcal T_N $:
\[
\norm{f - S_Nf}_2 \leq \norm{f - \sigma_Nf}_2
\leq \norm{f - \sigma_Nf}_\infty \to 0
\]
(a desigualdade média porque a média de $\abs\cdot^2$ está em
mais o sup ao quadrado). Pitágoras $\norm f_2^2 = \sum_{\abs
n\leq N}\abs{c_n}^2 + \norm{f - S_Nf}_2^2$ então passa para o
limite: Parseval para cada \omterm{def:b2:metric:continuity}{contínuo} periódico $ f $.

\textbf{9.} Deixar $ t_1, \dots, t_p $ sejam os saltos de $ f $ em um
período, $ M = \norm f_\infty $. Para pequenos $\eta $, definir $ g = f $
fora dos intervalos $\intoo{t_j - \eta}{t_j + \eta}$ e por
o acorde afim em cada intervalo: $ g $ é \omterm{def:b2:metric:continuity}{contínuo},
periódico, $\norm g_\infty \leq M $, e
\[
\norm{f - g}_2^2 \leq \frac{1}{2\pi}\,p\cdot(2M)^2\cdot2\eta
\leq \varepsilon^2
\]
para $\eta $ pequeno. Bessel faz $ S_N $ uma contração para
$\norm\cdot_2$, então
\[
\norm{f - S_Nf}_2
\leq \norm{f - g}_2 + \norm{g - S_Ng}_2 + \norm{S_N(g -
f)}_2
\leq 2\varepsilon + \norm{g - S_Ng}_2 ,
\]
e a questão 8 dá $\limsup_N\norm{f - S_Nf}_2 \leq
2\varepsilon $ para cada $\varepsilon $: $\norm{f - S_Nf}_2 \to
0$, e Pitágoras produz Parseval para cada pedaço
\omterm{def:b2:metric:continuity}{contínuo} $ f $: o ``admitido'' em
\cref{thm:b2:fourier:parseval} agora é um teorema.

\textbf{10.} A onda quadrada tem $\norm f_\infty = 1$, então
pergunta 5 dá $\abs{\sigma_N(f)} \leq 1$ em todos os lugares e para
cada $ N $ --- sem overshoot, nunca --- enquanto
\cref{exo:b2:fourier:8} mostra $\sup_xS_{2N-1}(f)(x) \to
\approx 1.179$. A única propriedade responsável: $ F_N \geq 0$, então
$\sigma_Nf(x)$ é um peso \emph{média} de valores de $ f $
e nunca pode sair $\intcc{\min f}{\max f}$; $ D_N $ leva
valores negativos, então $ S_N $ pode.

\textbf{11.} $\abs{\sin N\theta} \leq N\abs{\sin\theta}$ por
indução ($\abs{\sin(N{+}1)\theta} \leq
\abs{\sin N\theta}\abs{\cos\theta} +
\abs{\cos N\theta}\abs{\sin\theta} \leq
(N+1)\abs{\sin\theta}$): com $\theta = \frac u2$,
\[
F_N(u) = \frac{\sin^2\frac{Nu}2}{N\sin^2\frac u2}
\leq \frac{N^2\sin^2\frac u2}{N\sin^2\frac u2} = N .
\]
Concavidade de $\sin $ sobre $\intcc{0}{\frac\pi2}$ dá
$\sin\frac u2 \geq \frac{u}{\pi}$ para $0 \leq u \leq \pi $, então
$ F_N(u) \leq \frac{1}{N(u/\pi)^2} = \frac{\pi^2}{Nu^2}$.

\textbf{12.} Pela uniformidade e divisão em $\frac1N $:
\[
\frac{1}{2\pi}\int_{-\pi}^{\pi}\abs uF_N
= \frac1\pi\int_0^\pi uF_N
\leq \frac1\pi\Bigl(\int_0^{1/N}uN\,\dd u +
\int_{1/N}^{\pi}\frac{\pi^2}{Nu}\,\dd u\Bigr)
= \frac{1}{2\pi N} + \frac{\pi\ln(\pi N)}{N} .
\]
Para $ N \geq 2$: $\ln(\pi N) \leq \bigl(1 +
\frac{\ln\pi}{\ln2}\bigr)\ln N \leq 2.66\ln N $ e
$\frac{1}{2\pi N} \leq \frac{\ln N}{N}$, então o momento é
$\leq \frac{9\ln N}{N}$.

\textbf{13.} Para $ L $-Lipschitz $ f $:
\[
\abs{\sigma_Nf(x) - f(x)}
\leq \frac{1}{2\pi}\int\abs{f(x+u) - f(x)}F_N(u)\dd u
\leq L\cdot\frac{1}{2\pi}\int\abs uF_N
\leq \frac{9L\ln N}{N},
\]
\omterm{def:b2:funcseq:def}{uniformemente} em $ x $.

\textbf{14.} $ c_0(e_1) = 0$ dá $ S_0(e_1) = 0$, enquanto
$ S_n(e_1) = e_1$ para $ n \geq 1$: $\sigma_N(e_1) =
\frac{N-1}{N}e_1$, então $\norm{\sigma_Ne_1 - e_1}_\infty =
\frac1N $. Mesmo para todo esse sinal com banda limitada, a taxa é
$\frac1N $: Fejérrrr satura, exatamente como o operador de Bernstein
satura em $\frac1n $ (Voronovskaia, no
problema de fim de semana do capítulo sobre sequências de funções).

\textbf{15.} Se $ f $ desaparece em $\intoo{x-\delta}{x+\delta}$,
então $\sigma_Nf(x) = \frac{1}{2\pi}\int_{\delta \leq \abs u
\leq \pi}f(x+u)F_N(u)\dd u $, de valor absoluto no máximo
$\frac{\norm f_\infty}{N\sin^2(\delta/2)} \to 0$: o Cesàro
significa em $ x $ só veja $ f $ aproximar $ x $.

\textbf{16.} Dado $\intcc ab $ e $\varepsilon $, escolher
\omterm{def:b2:metric:continuity}{contínuo} $1$-periódico afim por partes $\varphi^\pm $ com
$\varphi^- \leq \mathbf 1_{\intcc ab} \leq \varphi^+$ e
$\int_0^1(\varphi^+ - \varphi^-) \leq \varepsilon $ (trapézios
com inclinações em intervalos de comprimento total $\varepsilon $).
Então
\[
\limsup_N\frac{\#\{n \leq N : x_n \in \intcc ab\}}{N}
\leq \lim_N\frac1N\sum_{n\leq N}\varphi^+(x_n)
= \int_0^1\varphi^+ \leq b - a + \varepsilon ,
\]
e simetricamente $\liminf \geq b - a - \varepsilon $: o
proporção tende a $ b - a $.

\textbf{17.} Para $ f = \eu^{2\iu\pi k\cdot}$ com $ k \neq 0$ o
hipótese dá o limite $0 = \int_0^1f $; para $ k = 0$ ambos
lados são $1$; linearidade lida com todos os trigonométricos
polinômio. Para \omterm{def:b2:metric:continuity}{contínuo} $1$-periódico $ f $ e $\varepsilon >
0$, questão 6 (transportado por $ t = 2\pi x $) fornece um
polinômio trigonométrico $ P $ com $\norm{f - P}_\infty \leq
\varepsilon $:
\[
\Bigl|\frac1N\sum_{n\leq N}f(x_n) - \int_0^1f\Bigr|
\leq 2\varepsilon +
\Bigl|\frac1N\sum_{n\leq N}P(x_n) - \int_0^1P\Bigr|
\longrightarrow 2\varepsilon .
\]
Com a pergunta 16: $(x_n)$ é equidistribuído.

\textbf{18.} Para $ k \neq 0$ e irracional $\alpha $, $ w =
\eu^{2\iu\pi k\alpha} \neq 1$:
\[
\Bigl|\sum_{n=1}^{N}w^n\Bigr|
= \Bigl|\frac{w(w^N - 1)}{w - 1}\Bigr|
\leq \frac{2}{\abs{1 - w}} ,
\]
um limite independente de $ N $; dividindo por $ N $ dá a Weyl
critério, e a questão 17 conclui: $(\{n\alpha\})$ é
equidistribuído.

\textbf{19.} Cada subintervalo recebe proporção assintótica
igual ao seu comprimento, em particular um número infinito de pontos:
$(\{n\alpha\})$ é denso. A equidistribuição é mais forte: uma
sequência pode ser densa enquanto passa quase todo o seu tempo em um
canto (a densidade diz onde a sequência \emph{vai},
equidistribuição diz \emph{com que frequência}).

\textbf{20.} $2^n $ tem dígito inicial $1$ se $10^m \leq 2^n <
2\cdot10^m $ para alguns $ m $, ou seja,\ se $\{n\log_{10}2\} \in
\intco{0}{\log_{10}2}$. Irracionalidade: $\log_{10}2 =
\frac pq $ daria $2^q = 10^p = 2^p5^p $, impossível para $ p
\geq 1$ por fatoração única. Teorema de Weyl (questão 18
com $\alpha = \log_{10}2$; o meio aberto intervalo é apertado
entre os fechados de comprimentos próximos) dá a proporção
$\log_{10}2 \approx 0.301$: os primeiros dígitos de $2^n $ siga
Lei de Benford.

\textbf{21.} $ z $ é $ C^1$, então sua série de Fourier converge
\omterm{def:b2:funcseq:series}{normalmente} com soma $ z $ (\cref{thm:b2:fourier:parseval}~(1)),
e $ c_n(z') = \iu n\,c_n $. Desde $\abs{z'} = \frac{L}{2\pi}$
é constante,
\[
\int_0^{2\pi}\abs{z'}^2\dd t
= 2\pi\Bigl(\frac{L}{2\pi}\Bigr)^{\!2}
= \frac{L^2}{2\pi} ,
\]
e Parseval aplicado ao \omterm{def:b2:metric:continuity}{contínuo} $ z'$ dá
$\frac{1}{2\pi}\int_0^{2\pi}\abs{z'}^2 =
\sum_n\abs{\iu nc_n}^2$, ou seja\ $\frac{L^2}{2\pi} = 2\pi\sum_n
n^2\abs{c_n}^2$.

\textbf{22.} O Parseval polarizado $\frac{1}{2\pi}\int\conj
fg = \sum\conj{c_n(f)}c_n(g)$ segue de Parseval aplicado a
$ f + g $ e $ f + \iu g $ (identidade de \omterm{def:b2:quadratic:def}{polarização}), ambos
\omterm{def:b2:metric:continuity}{contínuo}. Com $ f = z $, $ g = z'$:
\[
A = \frac12\,\Im\int_0^{2\pi}\conj z\,z'
= \pi\,\Im\sum_n\conj{c_n}(\iu n c_n)
= \pi\sum_n n\abs{c_n}^2 .
\]

\textbf{23.} Combinando as questões 21 a 22:
\[
L^2 - 4\pi A
= 4\pi^2\sum_n n^2\abs{c_n}^2 -
4\pi^2\sum_n n\abs{c_n}^2
= 4\pi^2\sum_{n\in\Z}(n^2 - n)\abs{c_n}^2 \geq 0 ,
\]
desde $ n^2 - n = n(n-1) \geq 0$ para cada número inteiro. Igualdade
forças $ c_n = 0$ para todos $ n \notin \{0, 1\}$: $ z(t) = c_0 +
c_1\eu^{\iu t}$, um círculo de centro $ c_0$ e raio
$\abs{c_1} = \frac{L}{2\pi}$ (velocidade constante). A prova de Hurwitz
da desigualdade isoperimétrica: $ A \leq \frac{L^2}{4\pi}$,
apenas círculo.

\textbf{24.} Círculo de raio $ R $: $ L = 2\pi R $, $ A = \pi
R^2$: $ L^2 = 4\pi^2R^2 = 4\pi A $: igualdade. Quadrado do lado
$ a $: $ L^2 = 16a^2 > 4\pi a^2 = 4\pi A $ (como $16 > 4\pi \approx
12.57$). Para negativo $ n $, $ n^2 - n = n(n - 1)$ é um produto
de dois inteiros negativos: positivo --- então retrocedendo
modos custam área duas vezes. Velocidade constante inserida na questão 21,
convertendo $\int\abs{z'}^2$ em $\frac{L^2}{2\pi}$; por um
velocidade não constante, Cauchy--Schwarz dá $\int\abs{z'}^2 \geq
\frac{(\int\abs{z'})^2}{2\pi} = \frac{L^2}{2\pi}$, então o
a desigualdade sobrevive, com o círculo ainda sendo a única igualdade
caso.

\textbf{25.} (i) Tudo flui de $ F_N \geq 0$ (com unidade
média e concentração); $ D_N $ tem média unitária e concentração
de oscilação, mas não de positividade, e Gibbs é o preço.
(ii) Cesàro somabilidade da série de Fourier implica sua
Abel somabilidade com a mesma quantia (Frobenius, provado no
problema de fim de semana do capítulo da série de potência) --- o núcleo de Poisson
rota de \cref{exo:b2:fourier:12} é exatamente o método de Abel.
(iii) o teorema de Weyl precisava apenas de uma aproximação uniforme
(questão 6); a desigualdade isoperimétrica necessária Parseval
em si (questões 8, 21--22). (iv) O volume do Ano 3 prova
completude: as exponenciais formam uma base de Hilbert de $ L^2$,
Parseval torna-se uma isometria dos espaços de Hilbert, e Fejérr
teorema torna-se a afirmação de que esta isometria é computável
por médias positivas.
\end{solution}
